\section{Related Work}
\label{sec:related}

\para{Mixing laws for loss.}
\citet{ye2025datamixing} propose data mixing laws: validation loss is a predictable function of the mixture proportions, exponential in the proportions, which can be fitted on a few sampled mixtures and nested with step and model-size scaling laws to choose a mixture for a larger model.
\citet{shukor2025scaling} extend scaling laws to predict loss as a joint function of model size, token count, and domain weights.
Both target validation loss.
We use the same kind of small-run fitting, but we ask whether downstream accuracy follows similar shapes, and we compare the exponential form against a log-share form on held-out mixtures.

\para{Choosing mixtures with proxy models.}
RegMix \citep{liu2025regmix} trains many tiny models on random mixtures, fits a regression from mixture to a target loss, and transfers the best mixture to larger models.
It reports that web corpora, and not sources usually seen as higher quality such as Wikipedia, correlate most positively with downstream performance.
DoReMi \citep{xie2023doremi} and DoGE \citep{fan2024doge} learn domain weights with a proxy model.
These methods search for a good mixture.
Unlike them, we do not optimize a mixture; we characterize the shape of the dose--response curve along each domain and test a specific threshold for web text.
RegMix's finding on web text runs against the claim that heavy web text is harmful, which makes H4 a real test.

\para{Code in pre-training.}
\citet{aryabumi2025code} run ablations on models with 470M to 2.8B parameters and report relative gains of up to 8.2\% on natural-language reasoning and 4.2\% on world knowledge when code is added to text-only pre-training.
This result motivates H3.
It uses a few mixtures, so it does not show how the benefit changes with the amount of code.
We sweep code share over nine levels and measure its correlation with reasoning probes on random mixtures.

\para{Scaling laws and small-model evaluation.}
Scaling laws describe loss as a power law in data and parameters \citep{kaplan2020scaling,hoffmann2022chinchilla}.
At a fixed budget a domain's share is its data quantity, so a logarithmic dependence of in-domain loss on share is the natural prior for H1.
Work on small language models, such as the BabyLM challenge \citep{warstadt2023babylm}, evaluates mostly with likelihood-based and fine-tuning protocols because zero-shot benchmark accuracy is near chance at that scale.
We follow this practice and measure reasoning by fine-tuning transfer.
